\section{Conclusion}
\label{conclusions}

In this work, we presented Ascend KernelGen, a unified framework designed to automate the generation of high-performance kernels for NPU architectures. We identified that the primary barrier for LLMs in this domain is not merely syntax familiarity, but the lack of structured reasoning capabilities regarding hardware constraints, memory hierarchy management, and asynchronous pipeline synchronization.

To overcome this, we constructed Ascend-CoT, a domain-specific dataset that explicitly models the reasoning process of expert developers, and developed NPUKernelBench, a rigorous evaluation suite that goes beyond static code analysis to verify compilation, numerical correctness, and runtime performance on actual hardware. Our extensive experiments show that while general-purpose models struggle with the specificity of AscendC, our domain-adaptive training strategy—combining reasoning-oriented supervised fine-tuning with execution-guided reinforcement learning—yields substantial improvements. We demonstrated that our model can successfully generate complex operators that were previously out of reach, achieving competitive performance metrics.

While this study focuses on kernel-level correctness and basic optimization, several promising directions remain. First, we plan to extend our framework to better support Level-3 kernels, which remain challenging for current models. Second, we aim to integrate performance-aware reward models into the reinforcement learning stage to prioritize not just functional correctness, but also latency minimization and resource utilization close to the hardware roofline. Finally, we intend to explore the generalizability of our methodology to other emerging accelerator platforms, fostering a more inclusive ecosystem for AI-driven hardware software co-design.